\section{A complete conjugacy test}
Let $G$ be finite, let $B\subseteq G$ be a forbidden set, and let $N\leq G$. Our convention is $x^g=g^{-1}xg$.
\begin{lemma}\label{avoid}
The following conditions are equivalent:
\[
 N^g\cap B=\varnothing\quad(g\in G),\qquad
 N\cap\bigcup_{b\in B}\Cl_G(b)=\varnothing.
\]
For the second condition it suffices to test one representative of every conjugacy class of elements of $N$. The admissible subgroups are closed under conjugation and passage to subgroups.
\end{lemma}
\begin{proof}
An element $b$ lies in $N^g\cap B$ exactly when $b=x^g$ for some $x\in N$ and $b\in B$. This is precisely membership of $x$ in a $G$-class meeting $B$. Conjugation within $N$ does not change its $G$-class, proving the representative test. Both closure assertions follow directly from the first condition.
\end{proof}

Suppose that $G$ is the marked group of a regular polytope, that $K$ is the stabilizer of a facet, and that $\omega$ is a central involution of $K$. Assume the quotient under consideration has flags indexed by the cosets of $N$, and that every intersection $N^g\cap K$ is either $1$ or $\langle\omega\rangle$. These assumptions concern the actual quotient and local sections, not just an arbitrary permutation action.
\begin{proposition}\label{facets}
Put $m=|N\cap\Cl_G(\omega)|$ and $c=|\Cen_G(\omega)|$. If $h$ and $d$ count the facet orbits with stabilizer of order two and one, respectively, then
\[
 h=\frac{2cm}{|K||N|},\qquad
 d=\frac{|G|}{|K||N|}-\frac h2.
\]
Moreover the full automorphism group of the quotient is $N_G(N)/N$.
\end{proposition}
\begin{proof}
Facets are cosets of $K$. A coset $Kg$ has an order-two stabilizer in $N$ exactly when $\omega^g\in N$. There are $cm$ choices of $g$ satisfying this condition: each of the $m$ possible conjugates has $c$ preimages. Centrality of $\omega$ in $K$ makes this condition independent of the representative of $Kg$, so there are $cm/|K|$ such facets. Their $N$-orbits have size $|N|/2$, which gives the formula for $h$. Counting all $|G|/|K|$ facets by their orbit sizes gives the formula for $d$.

On flags the marked generators act by adjacency. A type-preserving automorphism of the quotient commutes with these generators and hence with their generated transitive action. An equivariant bijection of $N\backslash G$ is determined by the image of $N$. That image is $Ng$ with the same stabilizer precisely when $g$ normalizes $N$. The resulting permutations are $N_G(N)/N$. Conversely, these permutations preserve every flag adjacency and therefore the quotient faces and incidence. This proves the automorphism assertion.
\end{proof}

\begin{lemma}\label{maxima}
Let a conjugacy-invariant collection of subgroups of a finite group be closed under subgroups. Form a directed graph on its subgroup conjugacy classes by recording, for each representative, the ambient conjugacy classes of all its maximal proper subgroups. A vertex is maximal in the collection exactly when it has no incoming edge.
\end{lemma}
\begin{proof}
An incoming edge exhibits a proper containing member. Conversely, if $N$ is properly contained in a member, choose a containing member minimal above $N$. Finiteness makes $N$ a maximal subgroup of that member. Conjugating the containing member to the chosen representative gives the required recorded edge. Repeated edges are counted only once.
\end{proof}

